\scieloAbstractContainer{Resumo}{Nos últimos anos, o processo de metropolização tem sido articulado à acentuação das desigualdades
socioeconômicas e à consequente difusão de práticas securitárias sob a forma de muros, grades, câmeras
de segurança, entre outros. Assim, em diversas metrópoles latino-americanas a expansão urbana tem sido}{Palavras-chave:}{Segregação socioespacial. Securização}

\scieloAbstractContainer{Abstract}{\textbf{Introduction:} This paper presents a detection method for upper limb movement intention as part of a brain-
machine interfaces using EEG signals, whose final goal is to assist disabled or vulnerable people with activities
of daily living. \textbf{Methods:} EEG signals were recorded from six naïve healthy volunteers while performing a
motor task. Every volunteer remained in an acoustically isolated recording room. The robot was placed in front}{Keywords:}{Movement intention, Objective response detection, Statistical F test, Event-related desynchronization, Brain-machine interface.}
%{\noindent\fontsize{9}{10.8}\selectfont{Received: January 13, 2015; Accepted: April 19, 2015.}}

\begin{multicols}{2}
\scieloSectionContainer{Introduction}
\par{}Brain interfaces are commonly considered as the
ultimate means of interaction between humans and
machines or computers (Allison et al., 2012; Doud et al.,
2011; Hochberg et al., 2012; Sirvent Blasco et al.,
2012). The idea of controlling our environment using
brain activity opens the possibility for many potential
applications using these interfaces. For instance, a

\scieloSectionContainer{Conclusion}

\lipsum

\scieloSectionContainer{Materials and Methods}

\scieloSubSectionContainer{Animals}

\par{}Specimens representing 12 species of polyclads were
collected from different sites in the Caribbean. In addition,
a deep-sea specimen of \textit{Anocellidus profundus} Quiroga,
Bolaños and Litvaitis, 2006, from the North Pacific was
also included in the analysis (specimen courtesy of Dr.
Janet Voight, Field Museum of Natural History, Chicago,
Illinois, USA). For details of collection, fixation, and spe-
cies identification see Litvaitis \textit{et al}. (2010).

\lipsum
\scieloSectionContainer{Considerações finais}

\par{}O processo de metropolização tem sido acompanhado
de um importante intensificação das desigualdades sociais, potencializando a difusão de práticas securitárias.
A partir da percepção de um aumento da violência nas
grandes metrópoles, a segurança é uma mercadoria
capaz de ser adquirida e consumida. Nas periferias
metropolitanas, torna-se cada vez mais comum a
disseminação de empreendimentos exclusivos cujo
acesso é controlado através de tecnologias de segurança


\lipsum

\scieloImageContainerOneCol{0.5\textwidth}{img1.jpg}{meu titulo}

%\begin{center}
%{
%\vspace{8mm}
%\centerline{
%\includegraphics[width=\maxwidth{0.5\textwidth}]{img1.jpg}
%}
%\vspace{8mm}
%}
%\end{center}

\lipsum

\scieloImageContainerTwoCol{img2.jpg}{\textbf{Figure 1:} A specimen of each species was embedded in paraffin in its entirety, longitudinally sectioned}

\lipsum
\lipsum

\scieloImageContainerTwoCol{img2.jpg}{\textbf{Figure 1:} A specimen of each species was embedded in paraffin in its entirety, longitudinally sectioned}

\lipsum
\lipsum
\lipsum[1]

\begin{scieloReferencesContainer}[Referências]



\scieloReferencesItem{Billard, G., Chevalier, J., \& Madoré, F. (2005). \textit{Ville fermée, ville
surveillée: la sécurisation des espaces résidentiels en France
et en Amérique du Nord}. Rennes: Presses Universitaires
de Rennes}

\scieloReferencesItem{Caldeira, T. (2000). Cidade de muros. \textit{Crime, segregação e cidadania em São Paulo}. São Paulo\: EDUSP.}

\scieloReferencesItem{Campos, P., \& Mendonça, J. (2013). Estrutura sócio-espacial
e produção habitacional na Região Metropolitana de
Belo Horizonte: novas tendências. In Cardoso, A. (Ed),
\textit{O programa Minha Casa Minha Vida e seus efeitos territoriais}
(p. 67-91). Rio de Janeiro: Letra Capital.}

\scieloReferencesItem{King GJ, Keeling SD, Hocevar RA, Wheeler TT. The timing of treatment for Class II malocclusions in children: a literature review. Angle
Orthod. 1990;60(2):87-97. PMid:2111647.}

\scieloReferencesItem{Proffit WR. The timing of early treatment: an overview. Am J Orthod Dentofacial Orthop. 2006 Apr;129(4 Suppl):S47-9. http://dx.doi.
org/10.1016/j.ajodo.2005.09.014. PMid:16644417.}

\scieloReferencesItem{Franchi L, Baccetti T, De Toffol L, Polimeni A, Cozza P. Phases of the dentition for the assessment of skeletal maturity: a diagnostic performance
study. Am J Orthod Dentofacial Orthop. 2008 Mar;133(3):395-400. PMid: 18331939.}

\scieloReferencesItem{Bullock TH and Horridge GA (1965) Structure
and Function in the Nervous Systems of Invertebrates.
Vol. 2. W.H. Freeman and Co, San
Francisco, 1719 pp.}

\scieloReferencesItem{Carranza S, Baguñà J and Riutort M (1997)
Are the Platyhelminthes a monophyletic primitive
group? An assessment using 18S rDNA
sequences. Mol Biol Evol 14:485-497.}

\scieloReferencesItem{Cebrià F (2008) Organization of the nervous system in the model
planarian Schmidtea mediterranea: An immunocytochemical study. Neurosci Res 61:375-384.}

\scieloReferencesItem{Chien P and Koopowitz H (1972) The ultrastructure of neuromuscular systems in Notoplana acticola, a free-living polyclad flatworm. Z Zellforsch Mikroskop Anat 133:277-288.}


\scieloReferencesItem{Chien PK and Koopowitz H (1977) Ultrastructure of nerve plexus
in flatworms. III. The infra-epithelial nervous system. Cell
Tissue Res 176:335-347.}

\scieloReferencesItem{Day TA, Maule AG, Shaw C and Pax RA (1997) Structure-activity relationships of FMRFamide-related peptides contracting Schistosoma mansoni muscle. Peptides 18:917-921.}

\scieloReferencesItem{Egger B, Gschwentner R and Rieger R (2007) Free-living
flatworms under the knife: Past and present. Dev Genes Evol
217:89-104.}

\scieloReferencesItem{Egger B, Lapraz F, Tomiczek B, Müller S, Dessimoz C, Girstmair
J, Skunca N, Rawlinson KA, Cameron CB, Beli E et al.
(2015) A transcriptomic-phylogenomic analysis of the evolutionary relationships of flatworms. Curr Biol 25:1-7.}

\scieloReferencesItem{Ehlers U (1985) Das Phylogenetische System der Plathelminthes.
Gustav Fischer Verlag, Stuttgart, 317 pp.}

\scieloReferencesItem{Fernandes MC, Alvares EP, Gama P and Silveira M (2003) Serotonin in the nervous system of the head region of the land planarian Bipalium kewense. Tissue Cell 35:479-486.}

\scieloReferencesItem{Forest DL and Lindsay SM (2008) Observations of serotonin and
FMRFamide-like immunoreactivity in palp sensory structures and the anterior nervous system of spionid polychaetes. J Morphol 269:544-551.}

\scieloReferencesItem{Girstmair J, Schnegg R, Telford MJ and Egger B (2014) Cellular
dynamics during regeneration of the flatworm Monocelis sp. (Proseriata, Platyhelminthes). Evo Devo 5:e37.}

\scieloReferencesItem{Golding DM (1992) Polychaeta: Nervous system. In: Harrison
FW and Gardiner SL (eds) Microscopic Anatomy of Invertebrates. Vol. 7. Wiley-Liss, New York, pp 155-179.}

\scieloReferencesItem{Gustafsson MKS, Halton DW, Kreshchenko ND, Movsessian SO,
Raikova OI, Reuter M and Terenina NB (2002) Neuropeptides in flatworms. Peptides 23:2053-2061.}

\scieloReferencesItem{Hadenfeldt D (1929) Das Nervensystem von Stylochoplana
maculata und Notoplana atomata. Z wiss Zool 133:586638.}

\scieloReferencesItem{Halton DW and Gustafsson MKS (1996) Functional morphology
of the platyhelminth nervous system. Parasitology
113:S47-S72.}

\end{scieloReferencesContainer}

\scieloLicenseContainer{License information: This is an open-access article distributed under the terms of the Creative Commons Attribution License, which permits unrestricted use, distribution, and reproduction in any medium, provided the original work is properly cited.}


\end{multicols}
